The discussion about the fairness of standardized exams in university admissions is ongoing. Opponents of standardized tests suggest these exams remain inaccessible to marginalized groups, while proponents argue that alternative assessments are prone to manipulation and often reflect a family's income. In response, some universities have adopted test-optional policies, allowing the submission of applications without mandatory standardized test results. Nonetheless, how these policies affect the diversity of student bodies remains unclear (Gooch et al. (2024)). For instance, research by Belasco et al. (2015) shows that test-optional policies do not boost diversity in liberal arts colleges but do lead to increased application competition. Conversely, Bennett (2018) reports that private institutions experience greater diversity with test-optional admissions. Despite an expanding body of research, it is still uncertain which student groups benefit most from these policies. Moreover, further comprehensive studies are required to evaluate improvements in application outcomes across different student groups.

In 2020, amidst the COVID-19 pandemic, the UK's teacher-assigned grading temporarily supplanted standardized exams, inadvertently elevating the final grades for most students. The British government annulled the A-level exam—an in-person standardized test evaluated by external assessors—just two months before it was scheduled to occur. This decision led to confusion among universities and students, who had previously established grade requirements for admissions. To compensate for the absence of final exam grades needed for university entrance, the government tasked schools with submitting subjective grades derived from students' academic records. Consequently, teachers often awarded overly generous grades, altering the alignment between students and their grades. Allegations arose against private secondary schools for inflating grades to enhance their pupils' placement statistics\footnote{Refer to "A-levels 2021: Rise three times more in private schools," available at \url{https://www.tes.com/magazine/news/secondary/levels-2021-rise-three-times-more-private-schools}. Accessed on December 18, 2024.}, while improvements in female students' grades compared to their male counterparts were also noted\footnote{See "Boys set to close gap with girls for top A-level grades," available at \url{https://www.thetimes.com/uk/education/article/a-level-results-2024-boys-expected-higher-girls-kc7sqkpqt}. Accessed on December 29, 2024.}. This situation illustrates how substituting a standardized exam can redistribute grade allocation among students and socio-economic groups.

This paper studies how switching from standardized grading to teacher-based grading impacts final grades and compares them between different student, exam, and school categories. Previous literature documents how attributes of schools and students grant grades disproportionate to the student's actual skill (Alesina et al. 2024, Lavy 2009). I study whether the increments in final grades are similar for university applicants. I hypothesize that the impact differs for higher education takers since the population of students has completed their mandatory education, and the consequences of high performance are more closely attached to labor market success. Furthermore, none of the previous papers compares the magnitude of the grade change between distinct student bodies. The universal switch in grading methods was universally implemented across all secondary schools in the UK, allowing me to observe the final assigned grade across all bodies for making comparisons. A comparison of marginal increases in grades between schools needs to control for the rich attributes of the student regarding their application, demographics, and attending schools. I study which schools are more likely to assign grades conditionally based on the characteristics of the student and their application content. 

I estimate the changes across final grades across student bodies using an interrupted time series model fitted with administrative data on university applications. I control for the student's previous scores in a standardized national exam mandated to all university applicants in the UK. I establish with a classical event study design that there are no pretrends on final grades after controlling for the student's previous exam scores. As the central part of the analysis, I decompose the grade increments across student's exam subjects, demographic characteristics, and their schools that provided them with the predicted grades. Since the final grade distribution is censored at the top, a Tobit model was used to model the grade assignment process during the pandemic. Additionally, to stabilize the estimates for each school, we estimate the school-level effects using a Bayesian hierarchy model. Recent heterogeneous treatment literature (Linero 2024) discusses the inefficiency of maximum likelihood estimates of multilevel models in finite samples. I use a generalized linear model and estimate the coefficient using Bayesian estimation. Then I quantify the determinants of the school-level estimates on school-level variables to understand the mechanisms of school-level grade inflation. 

%I estimate how a switch between standardized exams to a teacher-based-grading changes the overall grade distribution and its associated reshuffling of higher grade to different student bodies. My analysis focuses on the distributional changes across students' achievements in their past standardized exam, subject for the students final exam, student demographics, and the school attributes. I compare the achieved final grades between the cohort with TBG and those in previous years that took the standardized exam. I control for the students achivements in a standardized national exam that students must take when starting their university application education. The inclusion of the standardized controls allows me to compare across students with similar achievements in their previous exams, while not worrying about grading standards shifting across different year cohorts. 

%Firstly, I compare changes in the achieved final grades across achieved grades in the nationally standardized exam. I compare the magnitudes of the grade increase across students with different achivements in the GCSE exams. I test whether the new method changed the final grade distribution of disproportionately, whereas past literature documented teachers to provide more help to students at the threshold of not meeting their acceptance grade (Neal and Schanzenbach 2010). This also lets me quantify the extent of how much the grade standards were loosened. Secondly, I test whether subject content influenced the extent of the grade change, by comparing the increments in the final grade across different subjects categories. A more quantitative subject group (e.g. mathematics, physics, etc) may be easier for teachers to predict the final grades, since the students achivements are quantifiable. In contrast, more vocational subjects (drama studies, psychology, etc) may burden teachers with difficulties in providing a strict performance ranking across students, or provide an estimate of the students' future performance. Thirdly, I examine the increments across students in different demographic groups. Literature highlights the possibility of teacher rewarding students differentially based on the students' ethnic background (Alesina et al 2024). Finally, I test whether different types of schools may uniformly assign better grades to their pupils. Privately funded schools may react improve their pupils grade more than publicly funded schools, in order to appear to provide better education to their pupils (Brunello and Rocco, 2008). I also test whether the extent of the universal inflation is attributed to any other attributes of the schools, such as the inspection results or success in placing student to elite schools in previous years. I separate the factors influencing the reshuffling of grades into student independent factors and school level factors to fully characterize the distributional shift.

I find differential increments of final grades across school types and their pupil quality. I find that independent schools inflated grades by 0.1 standard deviation points compared to other public schools. I find an inverse U shape pattern between the share of students within a school with good scores in the previous national exams and the extent of their inflation. I observe a negative relationship between schools, with less than half of their students scoring good grades on their national exams. A school in the 50th percentile of the quality distribution assigned 0.5 standard deviation points higher grades than the school in the 0th percentile. A positive relationship is observed for schools with more than half of its students performing well. Additionally, I find grade inflation occurred more often for students taking non-quantitative subjects. Exams in facilitating subject categories were less inflated than non-facilitating subjects. Non-facilitating subjects had 9 percentage points less probability of obtaining an A or $A^*$ than facilitating subjects. However, facilitating subjects' exams were also inflated with their probability of receiving an A or $A^*$ increasing by 8 percentage points compared to previous years. I do not find differences in increments of assigned grades across student demographics after controlling for the student's subject choice.

My findings highlight how privately funded schools, such as independent schools, react to opportunities to improve the grades
of their students faster than other schools. In contrast to public schools, independent schools run by privately raised tuition paid by the parents of their students. Independent schools compete with each other by the number of students they place in better universities. Infer that independent schools took advantage of the loophole in grade assignment more quickly than other types of schools. 

My results also demonstrate how teacher-based grading rewards students who take subjects with more subjective discretion over their grading. Teachers may find difficulties in ranking performances of subjective subjects, due to their own preferences of content and the non-numerical responses of students. Teachers may also be reluctant to assign bad grades to students when they cannot strictly rank students from each other. The smaller increments in grades for quantitative subjects indicate that when teachers are aware of the students' ability with a well-defined metric, they are less likely to deviate from the grades a student would receive in a standardized exam. 



My results imply that the beneficiaries of non-standardized exam spreads across a diverse range of student bodies. Students in underrepresented demographic groups benefit from a teacher-based grading scheme since such subjects are difficult for teachers to assign bad grades. Their application success rate should improve because many underrepresented students take non-facilitating subjects. I infer students from advantageous backgrounds will not switch to non-facilitating subjects to increase their average grade, since many elite universities have strict subject requirements. There should be a reshuffling of university assignments between students from advantageous backgrounds who did not perform so well in their A-level with ambitious subject choice and students from less affluent backgrounds with less ambitious subject choice but with better overall grades. However, teacher-based grading is prone to manipulation, by more affluent schools having a stronger incentive to inflate the pupils' grades. 

\vspace{1em}


\textbf{Literature Review}  \ Chetty, Demming, and Friedman 2023 document the role of admission policies on the demographic composition of elite universities and their subsequent success. They predict admission policies using non-standardized tests that help affluent students enter elite universities, allowing them to own the financial resources they can deploy to obtain good certifications. However, their analysis is confined to elite universities and does not shed light on other student populations. 

Other papers differ between discretional grading and standardized grading (Cullen and Roback 2006, Neal and Schanzenbach 2010, Lavy 2009). They show the incentives behind entities in school institutions to support student groups based on systems that reward student performance. However, their results focus on students from a small number of institutions and fail to make comparisons between institutions with different characteristics. 

My paper highlights the role of classification methods on the distributional consequences in the higher education system.

The second section introduces the literature and further elaborates on my contribution. The third section describes the background of the UK higher education system and its reaction to the pandemic. The fourth section discusses the data. The fifth section discusses the empirical model. The sixth section shows the empirical results. The seventh section concludes and discusses policy implications.
